% ====================================================================
% [v1.0.21 Q2 경쟁 초안 — 창 q2f2] 다클래스 lattice gas 와 전하 보존 반전 (반영 전 초안 — 원본 무수정)
%   배치 전제: ch1_sec02b_part0.tex 의 \subsection{전기화학 연결 ...}(sec:sm-electro) 블록 끝
%   (fig:sm-occ figure 뒤) 과 \subsection{거시 열역학으로 ...}(sec:sm-macro) 사이에 삽입.
%   신규 라벨(기존 라벨·식 번호·수치 불변 전제): sec:sm-mc, eq:sm-mc-factor, eq:sm-mc-N,
%   eq:sm-mc-fix, eq:sm-mc-balance. 인용은 원장 V1 기존 키 hill1960·mcquarrie1976 만.
%   표기 주의: ch1_preamble 에는 \oc 매크로가 없으므로 $U_\mathrm{oc}$ 로 적는다(Ch2 의 $U_\oc$ 와 동일 렌더).
% ====================================================================

%% ================= A. 신설 소절 전문 =================
\subsection{다클래스 lattice gas --- 전하 보존 반전 $\sum_jQ_j\,\xi_{\eq,j}(U_\mathrm{oc},T)=Q\,\bar x$ 와 요동이 보증하는 유일근}\label{sec:sm-mc}
\S\ref{sec:sm-lattice} 가 한 종류의 자리 $M$ 개를 인수분해 $\Xi_M=\Xi_1^{\,M}$ 으로 닫고
\S\ref{sec:sm-electro} 가 $\mu$ 를 측정 전위에 결선했다. 실제 전극은 자리가 한 종류가 아니라 전이
$j=1,\dots,N_p$(\S\ref{sec:notation})가 겹친 계이므로, 이 소절은 같은 인수분해를 다클래스로 넓혀
\emph{관측 평형 전위를 결정하는 전하 보존 음함수}가 대정준 구조의 제약 반전 그 자체임을 보인다 --- 본론
\S\ref{sec:eqpeak} 의 전하 보존식과 Chapter 2 의 겹침 유도가 받는 출발점이 여기서 놓이고, 반전 근의
유일성까지 \S\ref{sec:sm-site} 의 요동 언어로 닫힌다.

\textbf{(a) 출발 --- 전이 $j$ $=$ 자리 클래스 $M_j$.} \S\ref{sec:sm-lattice} 의 결과는 동등 자리 $M$ 개의
대정준 합이 단일 자리 인수의 곱(식~\eqref{eq:sm-factor})으로 갈라져 $\langle N\rangle=M\theta$ 를 준다는
것이었다. 흑연에서는 전이 $j$ 마다 일하는 자리들이 서로 같지 않다 --- stage 갤러리마다 삽입의 자유에너지
비용이 다르다. 그래서 전이 $j$ 를 \emph{자리 클래스}(site class) 하나로 읽는다: 클래스 $j$ 는 동등 자리
$M_j$ 개와 유효 자리 자유에너지 $\tilde\varepsilon_j(T)$(식~\eqref{eq:sm-epstilde} 의 클래스판 ---
\S\ref{sec:sm-electro} 의 결선에서 전이 중심 $U_j$ 로 환산될 값)를 갖고, 저장조 $\mu$ 는 전해질이
공통이므로 모든 클래스가 공유한다. 자리 종류가 여럿인 국소화 흡착의 이 설정은 표준 전개다 \cite{hill1960}.

\textbf{(b) 연산 --- 클래스 곱 인수분해.} 가정 2(자리 독립)를 클래스 구조 위에 다시 놓으면 --- 클래스 안의
자리들은 동등$\cdot$독립, 클래스 사이도 독립 --- 대정준 합이 전 자리에 걸친 곱으로 인수분해된다:
\begin{equation}
\Xi\;=\;\prod_{j=1}^{N_p}\Xi_{1,j}^{\,M_j},\qquad
\ln\Xi\;=\;\sum_{j=1}^{N_p}M_j\,\ln\Xi_{1,j},\qquad
\Xi_{1,j}\;=\;1+e^{-\beta(\tilde\varepsilon_j-\mu)}
\label{eq:sm-mc-factor}
\end{equation}
--- 식~\eqref{eq:sm-factor} 를 클래스마다 한 번씩 쓴 것이다. 두 전제의 경계를 명시한다. (i) $\Xi_{1,j}$ 의
닫힌 두-항 꼴은 클래스 내 상호작용이 없는 이상 극한($\Omega_j=0$)의 것이고, $\Omega_j\ne0$ 이면 각 자리의
이웃 점유를 클래스 평균 $\theta_j$ 로 대체하는 \S\ref{sec:sm-mf} 의 평균장, 곧 \emph{자기무모순장}
(self-consistent field) 한정으로만 같은 곱꼴이 유지된다 --- $\tilde\varepsilon_j$ 가 $\theta_j$ 에 의존하는
유효값으로 옮겨 가고, $\theta_j$ 는 그 자기무모순 조건(식~\eqref{eq:mu})의 근이다. (ii) 클래스 \emph{사이}의
독립은 근사다 --- staging 에서 한 갤러리의 채움이 이웃 stage 의 문턱을 바꾸는 클래스 간 상관은 이 곱에서
빠지며, 본론은 그 몫을 전이별 모수 $U_j\cdot\Omega_j\cdot w_j$ 의 \emph{값} 안에 흡수한다.

\textbf{(c) 중간식 --- 평균 입자수 $=$ 자리수 가중 점유합과 함량 고정.} 식~\eqref{eq:sm-gc} 의 셋째 식을
\eqref{eq:sm-mc-factor} 의 로그에 적용하면 항마다 \S\ref{sec:sm-site} 의 단일 자리 결과가 되살아난다:
\begin{equation}
\langle N\rangle\;=\;k_BT\,\frac{\partial\ln\Xi}{\partial\mu}
\;=\;\sum_{j=1}^{N_p}M_j\,\theta_j(\mu),\qquad
\theta_j\;=\;\frac{1}{1+e^{+\beta(\tilde\varepsilon_j-\mu)}}.
\label{eq:sm-mc-N}
\end{equation}
대정준에서는 $\mu$ 가 독립변수이고 $\langle N\rangle(\mu)$ 이 따라 나온다. 그러나 측정은 반대로 통과 전하를
고정한다 --- 전극에 남은 Li 함량을 $x$, 추출(탈리튬화) 분율을 $\bar x\equiv1-x$ 라 하면($\bar x$ ---
Chapter 2 의 식 (eq:implicit)과 같은 배향의 조성 좌표) 고정되는 것은 $\langle N\rangle$ 쪽이고 미지수가
$\mu$ 다:
\begin{equation}
\sum_jM_j\,\theta_j(\mu)\;=\;(1-\bar x)\sum_jM_j
\quad\Longleftrightarrow\quad
\sum_jM_j\,\big[1-\theta_j(\mu)\big]\;=\;\bar x\,\sum_jM_j.
\label{eq:sm-mc-fix}
\end{equation}
여기에 \S\ref{sec:notation} 의 두 환산 --- 진행률 $\xi_j=1-\theta_j$, 자리수의 전하 환산
$Q_j=(F/N_A)\,M_j$($M_j$ 개 자리가 다 비워질 때 흐르는 전하가 전이 $j$ 의 용량이다) --- 을 얹고,
\S\ref{sec:sm-electro} 의 결선~\eqref{eq:sm-eqcond} 를 클래스마다 적으면
($U_j\equiv(\mu_\mathrm{Li}^\mathrm{metal}-N_A\tilde\varepsilon_j)/F$) 지수가 $\beta(\tilde\varepsilon_j-\mu)
\to sF(V-U_j)/RT$ 로 옷을 갈아입어 $1-\theta_j$ 는 식~\eqref{eq:sm-logistic} 의
$\xi_{\eq,j}(V,T)=\{1+e^{-sF(V-U_j)/RT}\}^{-1}$ 그대로가 된다.

\textbf{(d) 박스 --- 전하 보존 음함수의 대정준 정체.} 전류 없이 이완해 읽는 평형(개방 회로) 전위를
$U_\mathrm{oc}$ 라 하면($V=U_\mathrm{oc}$ 에서 평가 --- Chapter 2 와 같은 기호) 식~\eqref{eq:sm-mc-fix} 는
\begin{equation}
\boxed{\;\sum_{j=1}^{N_p}Q_j\,\xi_{\eq,j}(U_\mathrm{oc},T)\;=\;Q\,\bar x,
\qquad Q_j=\frac{F}{N_A}\,M_j,\quad Q=\sum_jQ_j\qquad(s=+1)\;}
\label{eq:sm-mc-balance}
\end{equation}
로 닫힌다 --- Chapter 2 의 식 (eq:implicit)이 겹침 유도(그 문건 파생 A)의 출발점으로 세우는 전하 보존
음함수와 문자 그대로 같은 식이고, 본론 \S\ref{sec:eqpeak} 이 관측 $\dd Q/\dd V$ 를 닫을 때 미분하는 전하
보존식 $Q_\cell q=Q_\bg(V_n)+\sum_jQ_j\xi_j$ 의 전이 몫의 통계역학적 기원이 바로 이것이다(배경 몫
$Q_\bg$ 는 삽입 자리 밖 채널이라 클래스 셈 바깥에서 더해진다). 곧 전하 보존식은 피팅 편의로 놓은 규칙이
아니라 \emph{고정 $\langle N\rangle$ 에서 $\mu$(전위)를 미지수로 돌린 대정준 제약 반전}이며,
$U_\mathrm{oc}$ 가 식 안에 음함수로 앉는 것은 논리 공백이 아니라 정의상 그런 서식(implicit formulation)이다
--- 그 반전의 근이 실제로 하나뿐인 이유가 아래 검산 (i)이다.
\begin{verifybox}
(i) \emph{유일근의 통계역학 증명.} 식~\eqref{eq:sm-gc} 의 $\langle N\rangle=k_BT\,\partial\ln\Xi/\partial\mu$
를 $\mu$ 로 한 번 더 미분하면 $\partial\langle N\rangle/\partial\mu=\beta\,(\langle N^2\rangle-\langle
N\rangle^2)=\beta\,\mathrm{var}(N)\ge0$ --- 앙상블 일반 항등식이다 \cite{mcquarrie1976}. 독립
클래스~\eqref{eq:sm-mc-N} 에서는 분산이 자리별로 더해져 $\mathrm{var}(N)=\sum_jM_j\,\theta_j(1-\theta_j)>0$
(유한 $T$ 에서 $0<\theta_j<1$)이니, \emph{입자수 요동의 양성}이 $\langle N\rangle(\mu)$ 의 순단조를 만든다
--- 단일 자리 요동--응답~\eqref{eq:sm-flucres} $\mathrm{var}(n)=\theta(1-\theta)$ 를 자리$\cdot$클래스 합으로
올린 거시판이다. 결선 $\mu_\mathrm{Li}=\mu_\mathrm{Li}^\mathrm{metal}-sFV$(식~\eqref{eq:sm-eqcond} 의
재배열)가 단조 직선이므로 박스 좌변은 $U_\mathrm{oc}$ 에 연속 순증($V\!\uparrow\Rightarrow\mu\!\downarrow
\Rightarrow\theta_j\!\downarrow\Rightarrow\xi_j\!\uparrow$ --- 전위를 올리면 탈리튬화, signbox 규약과 일치)
이고 $U_\mathrm{oc}\to\mp\infty$ 에서 $0/Q$ 로 가므로, $\bar x\in(0,1)$ 마다 근이 존재하고 유일하다 ---
Chapter 2 계산 예제$\cdot$부록 B 솔버가 서술한 ``좌변이 단조라 유일근''의 증명이 이 한 줄이다. 단 이 단조는
닫힌 logistic 서식($\Omega_j=0$ 엄밀 --- \S\ref{sec:sm-electro} 의 (iii))의 것이다 --- $\Omega_j>2RT$ 의
균질 평균장 가지는 비단조(그림~\ref{fig:sm-mu})이고 그 창의 분기 구성은 \S\ref{sec:hys} 소관이며, 본론과
Chapter 2 가 실제로 푸는 전이별 logistic 서식(폭 $w_j$ 의 이중지위는 \S\ref{sec:width})에서는
단조$\cdot$유일근이 그대로 유지된다. (ii) \emph{확장 off --- $N_p=1$ 회수.} 클래스가 하나면
\eqref{eq:sm-mc-factor} 는 $\Xi=\Xi_1^{\,M}$(식~\eqref{eq:sm-factor})로, \eqref{eq:sm-mc-N} 은 $\langle
N\rangle=M\theta$ 로 \S\ref{sec:sm-lattice} 를 정확히 회수하고, 박스는 $\xi_\eq(U_\mathrm{oc},T)=\bar x$ 가
되어 닫힌 꼴로 뒤집힌다 --- 그 역함수가 다음 소절 Nernst 식~\eqref{eq:sm-nernst} 이고, $N_p\ge2$ 겹침에서는
닫힌 역이 없어 수치 유일근만 남는다. (iii) \emph{반전 읽기.} 대정준의 순방향은 $\mu\mapsto\langle N\rangle$,
측정의 순방향은 $\bar x\mapsto U_\mathrm{oc}$ 다 --- 결선이 $\mu$ 와 $V$ 를 일대일로 묶으므로 박스를
$U_\mathrm{oc}$ 에 대해 푸는 것은 고정 $\langle N\rangle$ 에 맞는 $\mu$ 를 찾는 것과 같은 일이고, ``전하
보존이 전위를 결정한다''의 통계역학 문장이 ``고정 $\langle N\rangle$ 이 $\mu$ 를 결정한다''이다.
\end{verifybox}

이로써 자리 하나의 점유에서 곡선 전체의 전위 결정 구조까지가 한 앙상블 언어로 이어졌다. 남은 다리는 이
통계역학 $\mu$ 가 본론의 거시 언어 $G\cdot H\cdot S$ 와 같은 양이라는 확인이며, 다음 소절이 그 다리와
Nernst 식으로 Part 0 을 마감한다.

% 자산(v1.0.21 반영 시 부여): [V21-nnn 다클래스 대정준 인수분해·전하 보존 반전·요동 유일근 소절 sec:sm-mc]

%% ================= B. keybox "Part 0 사다리" 추가 1항 (ch1_sec02b_part0.tex 기존 keybox 수정안) =================
%% 기존: ... 전위 손잡이~\eqref{eq:sm-eqcond} $\to$ logistic~\eqref{eq:sm-logistic}·Nernst~\eqref{eq:sm-nernst}.
%% 신안: ... 전위 손잡이~\eqref{eq:sm-eqcond} $\to$ logistic~\eqref{eq:sm-logistic} $\to$ 다클래스
%%       $\langle N\rangle{=}\sum_jM_j\theta_j$·전하 보존 반전~\eqref{eq:sm-mc-balance} $\to$ Nernst~\eqref{eq:sm-nernst}.
%% (사다리 순서 = 소절 배치 순서. 나머지 keybox 문구 불변.)

%% ================= C. 연결 참조 문장 (각 1문장 — 재유도 없이 참조만; Ch2 는 별도 컴파일이라 서술형 참조) =================
%% C-1. ch2_sec05_mixing.tex 파생 A — 식 eq:implicit 도입 문단 끝("... $Q=\sum_j Q_j$ 다." 뒤)에 삽입:
%%   이 음함수는 피팅 관례가 아니라 다클래스 lattice gas 의 대정준 평균 $\langle N\rangle=\sum_jM_j\theta_j$ 를
%%   고정 함량에서 $U_\oc$ 에 대해 뒤집은 제약 반전이며, 그 유도와 유일근의 요동 증명은 Chapter 1 Part 0
%%   (그 문건 식 (eq:sm-mc-balance))이 닫는다.
%% C-2. ch2_appB_codemap.tex B.1 — \code{solve\_U\_oc} 진입점 서술 문장 뒤에 삽입:
%%   \code{solve\_U\_oc} 가 전제하는 유일근은 구현의 요행이 아니라 통계역학의 보증이다 ---
%%   $\partial\langle N\rangle/\partial\mu=\beta\,\mathrm{var}(N)>0$(입자수 요동의 양성)이 음함수~\eqref{eq:implicit}
%%   좌변의 $U_\oc$ 단조를 만든다는 증명을 Chapter 1 Part 0(그 문건 식 (eq:sm-mc-balance))이 상술한다.
